\documentclass[a4paper,11pt]{article}

\usepackage{../../../utils/template}

\newgeometry{ left=1.0cm, right=1.0cm, top=0.8cm, bottom=0.8cm }

\pagestyle{empty}

% ===== 题目集版式宏 =====
% 章标题：独立成页大标题
\newcommand{\chaptermark}[1]{%
    \newpage
    \begin{center}
        {\Large\textbf{#1}}
    \end{center}
    \vspace{1em}}
% 节标题
\newcommand{\booksec}[1]{\vspace{0.55em}\noindent\textbf{#1}\vspace{0.2em}\par}
% 习题块标题（如「习题 1-1」）
\newcommand{\exgroup}[1]{\vspace{0.6em}\noindent\textbf{#1}\par\vspace{0.2em}}
% 例题：\examp{例 x.x}{题干}
\newcommand{\examp}[2]{\vspace{0.65em}\noindent\textbf{#1}\quad #2\par\vspace{0.3em}}
% 习题列表
\newenvironment{probs}
    {\begin{enumerate}[leftmargin=2.2em, itemsep=0.55em]}
    {\end{enumerate}}

% ===== 选项排版（均靠左）=====
% opts1：每个选项独占一行（长选项）
\newenvironment{opts1}
    {\par\vspace{0.2em}\noindent\renewcommand{\arraystretch}{1.25}\begin{tabular}{@{}p{0.96\linewidth}@{}}}
    {\end{tabular}\par\vspace{0.2em}}
% opts2：两个选项一行，各占 50%
\newenvironment{opts2}
    {\par\vspace{0.2em}\noindent\renewcommand{\arraystretch}{1.25}\begin{tabular}{@{}p{0.475\linewidth}p{0.475\linewidth}@{}}}
    {\end{tabular}\par\vspace{0.2em}}
% opts4：四个选项一行，各占 25%
\newenvironment{opts4}
    {\par\vspace{0.2em}\noindent\renewcommand{\arraystretch}{1.25}\begin{tabular}{@{}p{0.235\linewidth}p{0.235\linewidth}p{0.235\linewidth}p{0.235\linewidth}@{}}}
    {\end{tabular}\par\vspace{0.2em}}

\begin{document}


% frag1: p01-p05
% 说明：p01–p05 为张宇《高等数学18讲》第 1 章（函数 极限 连续）第 2 章前内容页，
% 均为「无穷小/极限计算」一节的内页，页面上未出现章标题或节标题原文，故未写 \chaptermark/\sectionmark。
% 只转录题目（例 1.1–例 1.5），【注】【分析】【解】及理论正文一律截断不录。

% ---- p01 ----
% 书籍第 2 页：u^v 未定式整体判定的【注】之后的例 1.1（填空）。
\examp{例1.1}{$\lim\limits_{x\to 0^{+}}\dfrac{x^{x}-1}{\ln x\cdot\ln(1-x)}=\underline{\hspace{3em}}$.}

% ---- p02 ----
% 例 1.2：变上限积分参与的未定式计算。其后【分析】【解】【注】截断；
% 页末「2. 未定式局部判定（D1（常规操作）+D23（化归经典…）」为理论小节标题，无题目。
\examp{例1.2}{计算 $\lim\limits_{x\to 0}\left[\dfrac{1+\displaystyle\int_{0}^{x}\dfrac{\sin t}{t}\,\mathrm{d}t}{x}-\dfrac{1}{\ln(1+x)}\right]$.}

% ---- p03 ----
% 侧排页：例 1.3（选择题），间断点个数判定。其后【解】截断。
\examp{例1.3}{设 $f(x)=\begin{cases}\lim\limits_{t\to x}\left(\dfrac{x-1}{t-1}\right)^{\frac{t}{x-t}}, & x\neq 1,\\[1ex] 0, & x=1,\end{cases}$ 则 $f(x)$ 的第二类间断点的个数为（\quad）.}
\begin{opts4}
(A) $0$ & (B) $1$ & (C) $3$ & (D) $\infty$
\end{opts4}

% ---- p04 ----
% 侧排页：例 1.4（填空），等价无穷小比较。其后【解】及正文「(4) 变上限积分型」等理论内容截断。
\examp{例1.4}{当 $x\to 0$ 时，$\cos\left(e^{\frac{x^{2}}{2}}-1\right)-1\sim cx^{k}$，则 $ck=\underline{\hspace{3em}}$.}

% ---- p05 ----
% 侧排页：例 1.5（选择题），高阶无穷小运算律。其后【解】截断；
% 页面其余为理论（(6) 带头大哥型、「4. 常用的无穷大量的比较」及【注】内的举例），均非编号题目，不录。
\examp{例1.5}{设函数 $f(x)$，$g(x)$ 在 $x=0$ 的某去心邻域内有定义且恒不为 $0$，若当 $x\to 0$ 时，$f(x)$ 是 $g(x)$ 的高阶无穷小，则当 $x\to 0$ 时，有（\quad）}
\begin{opts2}
(A) $f(x)+g(x)=o(g(x))$ & (B) $f(x)g(x)=o(f^{2}(x))$ \\
(C) $f(x)=o(e^{g(x)}-1)$ & (D) $f(x)=o(g^{2}(x))$
\end{opts2}

% frag2: p06-p10

% ---- p06 ----
% 第 6 页：顶部为上一页解答的延续（无穷大阶的比较：当 x→+∞ 时 ln^p x ≪ x^a ≪ a^x；当 n→∞ 时 ln^p n ≪ n^a ≪ a^n ≪ n! ≪ n^n，
% 以及 lim ln^p n/n^q = 0、lim n^q/a^n = 0、lim a^n/n! = 0、lim n!/n^n = 0），属知识点总结，非题目。
\examp{例1.6}{已知函数 $f(x)=\dfrac{\int_{0}^{x}\ln(1+t^{2})\,\mathrm{d}t}{x^{\alpha}}$ 在 $(0,+\infty)$ 上有界，则 $\alpha$ 的取值范围是（\quad）.}
\begin{opts4}
(A) $-1<\alpha<3$ & (B) $1<\alpha\leqslant 3$ & (C) $0<\alpha\leqslant 3$ & (D) $\alpha>3$
\end{opts4}

% ---- p07 ----
% 第 7 页：顶部 ①② 为知识点总结（① x→∞ 时极限存在的充要条件；② lim 局部中 f−f 型函数用拉格朗日中值定理：
% f(b)−f(a)=f'(ξ)(b−a)），非题目。
\examp{例1.7}{$\lim\limits_{n\to\infty}n^{2}\bigl(\sqrt[n]{2}-\sqrt[n+1]{2}\bigr)=\underline{\hspace{5em}}$.}

% ---- p08 ----
% 第 8 页：顶部为知识点总结（化归经典形式示例：e^{sin x}−e^{tan x}=e^{tan x}(e^{sin x−tan x}−1)；
% e−(1+1/x)^x=e−e^{x ln(1+1/x)}=e[1−e^{x ln(1+1/x)−1}]）；中部 ★★★④ 见到计算 lim_{x→∞}[f(x)−ax] 时
% 将 f(x) 与 ax 通过恒等变形“合在一起”，均为知识点总结，非题目。
\examp{例1.8}{$\lim\limits_{x\to+\infty}x^{2}\left(3^{\frac{1}{x}}-3^{\frac{1}{x+1}}\right)=\underline{\hspace{5em}}$.}

\examp{例1.9}{$\lim\limits_{x\to\infty}\left[x\ln\left(\mathrm{e}+\dfrac{1}{x-1}\right)-x\right]=\underline{\hspace{5em}}$.}

% ---- p09 ----
% 第 9 页：顶部为例 1.9 解答的延续；左边缘为上一页知识点示例的延续（cos(sin x)−cos(tan x)，考虑如下解法）。
\examp{例1.10}{$\lim\limits_{x\to+\infty}\left[\dfrac{x^{1+x}}{(1+x)^{x}}-\dfrac{x}{\mathrm{e}}\right]=\underline{\hspace{5em}}$.}

% \booksec{二、判定连续与间断（$O_{2}$（盯住目标2））}

% \booksec{1. 常见备选点判定（$D_{1}$（常规操作））}

% 该小节下 ①（e^{1/x} ⇒ x=0 为无定义点，蓝字注：D_{21}（观察研究对象），盯住无定义点与分段点，
% 尤其是分母表达式的零点）为知识点总结，非题目。

% ---- p10 ----
% \booksec{3. 按定义作出结论（$D_{22}$（转换等价表述））}

% 该小节下 ①跳跃间断点、②可去间断点、③无穷间断点、④振荡间断点（附等价表述与小图）为知识点总结，非题目。
\examp{例1.11}{函数 $f(x)=|x|^{\frac{1}{(1-x)(x-2)}}$ 的第一类间断点的个数是（\quad）.}
\begin{opts4}
(A) $3$ & (B) $2$ & (C) $1$ & (D) $0$
\end{opts4}

% frag3: p11-p15
% 来源：考研教辅影印照片（横向拍摄，已旋转转正后转录）；本组 5 页均为例题页，
% 无「习题 x-x」整块；所有【证】【解】【注】及旁注、思维导图框均属解答/说明，未转录。

% ---- p11 ----
% 书页 10：页首为上一例题（选择题）解答的结尾（“对于 x=2，lim|x|^{1/((1-x)(x-2))}=e^{+∞}，故 x=2 是第_类间断点．综上，第一类间断点的个数是 1，选(C).”），属解答，未转录。
% \booksec{三、研究 $x\to\bullet$ 时 $f(x)$ 的微观性态（$O_{3}$（盯住目标 3））}
% \booksec{1. 定义法（$D_{22}$（转换等价表述））}
\examp{例 1.12}{设 $f(x)$ 具有一阶连续导数，且 $\lim\limits_{x\to+\infty}[x-f(x)]=0$，$f(1)<1$，证明：(1) 存在 $\xi>1$，使得 $|\xi-f(\xi)|<1-f(1)$；(2) 存在 $\eta>1$，使得 $f'(\eta)>1$.}

% ---- p12 ----
% 书页 11：页首蓝框为上一例题解答的结尾（“取 ε=…； (2) 反证法（此处只证 A≥0 的情况）……故可得 A≥0，证毕.”），
% 其后为例 1.13 的【证】及【注】框，均属解答/说明，未转录。
% 注意：本节小标题「2. …」未出现在本组照片中（书页 11 顶部被照片裁切），编号从 1 直接跳到 3。
\examp{例 1.13}{设 $f(x)$ 单调减少，$\lim\limits_{x\to+\infty}f(x)=0$，证明 $f(x)\geqslant 0$.}
% \booksec{3. 夹逼准则（$D_{1}$（常规操作）$+\ D_{22}$（转换等价表述）$+\ D_{23}$（化归经典形式））}

% ---- p13 ----
% 书页 12：页首为夹逼准则的定理条件「(1) h(x)⩽f(x)⩽g(x)；(2) lim g(x)=A，lim h(x)=A. 则 lim f(x) 存在，且 lim f(x)=A.」——定理陈述，非题目，未转录。
\examp{例 1.14}{设 $x\geqslant 0$，记 $x$ 到 $2k$ 的最小距离为 $f(x)$，$k=0,1,2,\cdots$.\\
(1) 证明 $f(x)$ 以 2 为周期，写出其在 $[0,2]$ 上的表达式并画出 $f(x)$ 的图像；\\
(2) 求 $\lim\limits_{x\to+\infty}\dfrac{\int_{0}^{x} f(t)\,\mathrm{d}t}{x}$.}

% ---- p14 ----
% 书页 13：页首「至于 x→−∞，x→+∞ 的情况，考生不难得出结论.」为例 1.14 解答的结尾，未转录。
\examp{例 1.15}{已知函数 $f(x)$ 的定义域是 $[0,+\infty)$，且满足 $f(0)=1$，$f'(x)=\dfrac{1}{f^{2}(x)+x^{2}}$，求证：$\lim\limits_{x\to+\infty}f(x)$ 存在，且 $\lim\limits_{x\to+\infty}f(x)\leqslant 1+\dfrac{\pi}{2}$.}

% ---- p15 ----
% 书页 14：跨入第 2 讲（数列极限）的开篇内容；左页为第 2 讲思维导图框，右页顶部为「压缩映射」说明文字及例 2.1 的【解】，
% 均属说明/解答，未转录；本页页面上未出现「第 2 讲」章标题文字（仅据例题编号 2.x 推断所属讲次）。
\examp{例 2.1}{若 $x_{1}=1$，$x_{n+1}=\sqrt{4+3x_{n}}$，$n=1,2,\cdots$，证明数列 $\{x_{n}\}$ 收敛，并求 $\lim\limits_{n\to\infty}x_{n}$.}
\examp{例 2.2}{若 $x_{n+1}=\dfrac{\pi}{2}+\dfrac{1}{2}\cos x_{n}$（$n=1,2,\cdots$），$x_{1}=\pi$，证明数列 $\{x_{n}\}$ 收敛，并求 $\lim\limits_{n\to\infty}x_{n}$.}

% frag4: p16-p20
% 说明：p16–p20 为与 frag1 同一本书（张宇《高等数学18讲》）第 2 讲（数列极限）的连续内容页
% （p17 左下角可见书页码 16），均为例题内页：无章标题、无习题块；仅 p19 出现两个小节标题（二、三）。
% 本批只转录题目（例 2.3–例 2.8，共 6 题）；【分析】【证】【解】、蓝色手写批注及理论方法条目
% （「(2) 若 x_{n+1}=f(x_n)…」「①…④」等）一律截断不录。照片均为横排 90° 侧拍，已旋转校读。

% ---- p16 ----
% 本页：顶部为例 2.2 解答结尾（「又当 n→∞ 时，(1/2)^n→0^+，故数列 {x_n} 收敛，且 lim x_n = π/2」）
% 与理论条目「(2) 若 x_{n+1}=f(x_n)，f(x) 不易求导或 f'(x) 不满足 |f'(x)|≤k<1.（D₁（常规操作））」
% 「①直接比较 x_{n+1} 与 x_n 的大小，定单调.」，均非题目。例 2.3 为本页唯一题目，其后【分析】【解】截断。
\examp{例2.3}{若 $0<x_n<1$，$x_{n+1}=1-\sqrt{1-x_n}$（$n=1,2,\cdots$），求：(1) $\lim\limits_{n\to\infty}x_n$；(2) $\lim\limits_{n\to\infty}\dfrac{x_{n+1}}{x_n}$.}

% ---- p17 ----
% 书页 16：顶部为例 2.3 第(2)问解答结尾（「(2) 由(1)可知，lim x_{n+1}/x_n = … = 1/2」）
% 与理论条目「②作差 x_{n+1}-x_n，根据正负定单调.」，非题目。本页两题（例 2.4、例 2.5），其后【解】截断。
\examp{例2.4}{设 $x_{n+1}=\dfrac{x_n^{2}}{2(x_n-1)}$（$n=1,2,\cdots$），$x_1>1$，证明数列 $\{x_n\}$ 收敛，并求 $\lim\limits_{n\to\infty}x_n$.}
\examp{例2.5}{设 $x_{n+1}=\sqrt{x_n(2-x_n)}$（$n=1,2,\cdots$），$0<x_1<2$，证明数列 $\{x_n\}$ 的极限存在，并求此极限.}

% ---- p18 ----
% 跨页照片：左缘重复书页 16（例 2.5 解答片段，已录，不重复）；右页（书页 17）顶部为例 2.5 解答结尾
% （「a = lim x_n = lim √(x_{n-1}(2-x_{n-1})) = √(a(2-a))，解得 a=1 或 a=0（舍去），故 lim x_n = 1」）
% 与理论条目「④根据题设提示(往往是第(1)问)，判有界或单调.」，非题目。本页一题（例 2.6），其后【证】截断。
\examp{例2.6}{(1) 证明对任意正整数 $n$，都有 $\dfrac{1}{n+1}<\ln\left(1+\dfrac{1}{n}\right)<\dfrac{1}{n}$ 成立； \\
(2) 设 $a_n=1+\dfrac{1}{2}+\cdots+\dfrac{1}{n}-\ln n$（$n=1,2,\cdots$），证明数列 $\{a_n\}$ 收敛.}

% ---- p19 ----
% 书页（约 18）：两个小节标题照录；标题下理论句「比较 x_n,x_{n+1} 的大小，定 x_n 的上、下界.（注：可能考虑 f(x) 的最值）」
% 与「对于 x_{n+1}=f(x_n)，若初值 x_1 仅给定取值范围，可能需要分情况讨论.」非题目。
% 例 2.7 其后【解】截断；例 2.8 为选择题，本页 (C)(D) 选项被手指遮挡、下缘裁切，文字据同页完整照片 p20 校补。
% \booksec{二、见到 $f(x_n,x_{n+1})$，且其是含 $x_n,x_{n+1}$ 的不等式关系（$D_1$（常规操作））}
\examp{例2.7}{若 $(1-x_{n+1})x_n>\dfrac{1}{4}$（$n=1,2,\cdots$），$0<x_n<1$，证明数列 $\{x_n\}$ 收敛，并求 $\lim\limits_{n\to\infty}x_n$.}
% \booksec{三、初值 $x_1$ 对 $x_{n+1}=f(x_n)$ 的敛散性的影响（$D_1$（常规操作））}
\examp{例2.8}{设数列 $\{a_n\}$ 满足 $\begin{cases}a_1>2,\\ a_{n+1}=a_n-\dfrac{a_n-3}{a_n-2},\end{cases}$ $n=1,2,3,\cdots$，则（\quad）.}
\begin{opts2}
(A) $\{a_n\}$ 收敛于大于 $3$ 的数 & (B) $\{a_n\}$ 收敛于 $3$ \\
(C) $\{a_n\}$ 发散 & (D) $\{a_n\}$ 的敛散性与 $a_1$ 有关
\end{opts2}

% ---- p20 ----
% 跨页照片：左页与 p19 为同一书页的完整版（例 2.8 及选项已录，不重复；其【解】应选(B) 起为解答，截断）。
% 右缘仅露出下一页（书页 19）的窄条：例 2.8 解答结尾（「设 a_1>2 且 a_1≠…」「显然当 a_1=3 时…」
% 「综上，{a_n} 收敛…」）、小节标题「四、双通…（D₁（常规…）」（标题被页缘裁断，仅余「四、双通」四字），
% 其下条目「(1) 将 a_n,b_n 满…小比阶；③常见放…」「(2) 令 c_n=a_n/b_n…」及题号「例2.9」——
% 题干均被裁去不可读，故无题可录，跳过（例 2.9 需下一张照片补全）。

% frag5: p21-p25
% 说明：p21-p25 为同一本书连续页（书页约 17-22），均属章「第 2 讲 数列极限」（章名取自 p24 页侧蓝色书眉）。
% 照片为书摊开拍摄，前后照片页面有重叠，重叠出现的题目只转录一次；只录题干，遇【解】截断。

% ---- p21 ----
% p21 上缘为上一页解答残句（“…界，若 a_n=f(n)”、“\frac{3}{2}\cdot\frac{1}{x-2}+1”、“…在 x\in I 上单调…”），非完整题目，不录。
% p21 另有前例解答的方法小结 (1)(2)（联立消元 / 令 c_n=\frac{a_n}{b_n} 转求极限），非题目，不录。
\examp{例 2.9}{设数列 $\{a_n\}$，$\{b_n\}$ 满足 $a_n+b_n=\dfrac{\pi}{2}$，$a_{n+1}=a_n+\cos a_n$，其中 $n=1,2,\cdots$，$a_1=1$，当 $n\to\infty$ 时，$b_n$ 与 $b_{n-1}^{a}$ 为同阶无穷小量，则 $a=(\quad)$.}
\begin{opts4}
(A) $1$ & (B) $2$ & (C) $3$ & (D) $5$
\end{opts4}
% 自【解】起为解答，按规则截断不录。

% ---- p22 ----
% p22 上部为例 2.9 解答的续页（“…得 A=A-\sin A，故 A=0，即 \lim b_n=0 … b_n=b_{n-1}-\sin b_{n-1}\sim\frac{1}{6}b_{n-1}^3，因此 a=3”），不录。
\examp{例 2.10}{设正项数列 $\{a_n\}$ 收敛于 $0$，若 $a_{n+1}=\cos b_n-\cos a_n$，$a_n\in\left(0,\dfrac{\pi}{2}\right)$，$b_n\in\left(0,\dfrac{\pi}{2}\right)$，且 $(1-b_n)^n=\cos b_n$，则 $\lim\limits_{n\to\infty}b_n^{\frac{\ln\cos b_n}{n}}=\underline{\hspace{3em}}$.}
% 自【解】起为解答，按规则截断不录。
% \booksec{五、复合函数的极限}
% 节下为「1.定理一(因变量极限定理)」的定理陈述（三条件），属理论非题目，不录。

% ---- p23 ----
% 例 2.11：题干跨 p23-p24 书页（p23 给函数设定与三张函数图，p24 顶部接选项行）。
\examp{例 2.11}{设 $u=g(x)=\begin{cases} x, & x \text{ 是有理数} \\ 0, & x \text{ 是无理数} \end{cases}$，$y=f(u)=\begin{cases} 1, & u\neq 0 \\ 0, & u=0 \end{cases}$，则 $\lim\limits_{x\to x_0}f[g(x)]$（\quad）.}
\begin{opts4}
(A) 等于 $0$ & (B) 等于 $1$ & (C) 等于 $x$ & (D) 不存在
\end{opts4}
% 自【解】起为解答，另有【注】与蓝色批注，均不录。

% ---- p24 ----
% p24 上部（书页 21 顶部）为上一题“则 \lim …”行及选项、解答、【注】，已并入上题，不重录。
% 节下小标题：2.定理二(中间变量极限定理)（定理陈述非题目，不录）。
\examp{例 2.12}{已知数列 $\{x_n\}$，其中 $-\dfrac{\pi}{2}\leqslant x_n\leqslant\dfrac{\pi}{2}$，则（\quad）.}
\begin{opts1}
(A) 若 $\lim\limits_{n\to\infty}\cos\sin x_n$ 存在，则 $\lim\limits_{n\to\infty}x_n$ 存在\\
(B) 若 $\lim\limits_{n\to\infty}\sin\cos x_n$ 存在，则 $\lim\limits_{n\to\infty}x_n$ 存在\\
(C) 若 $\lim\limits_{n\to\infty}\cos\sin x_n$ 存在，则 $\lim\limits_{n\to\infty}\sin x_n$ 存在，但 $\lim\limits_{n\to\infty}x_n$ 不一定存在\\
(D) 若 $\lim\limits_{n\to\infty}\sin\cos x_n$ 存在，则 $\lim\limits_{n\to\infty}\cos x_n$ 存在，但 $\lim\limits_{n\to\infty}x_n$ 不一定存在
\end{opts1}
% 自【解】起为解答，按规则截断不录。（例 2.12 在 p25 照片中又完整出现一次，系照片重叠，此处只录一次。）

% ---- p25 ----
% p25 上缘为前页残句（例 2.10 解答尾、例 2.11 批注、二维码），不录；例 2.12 重复出现，不重录。
\examp{例 2.13}{设正项数列 $\{a_n\}$ 单调增加，则以下选项中使得 $\{a_n\}$ 收敛的是（\quad）.}
\begin{opts2}
(A) $\left\{(1+a_n)^{\frac{1}{n}}\right\}$ 收敛于 $1$ & (B) $\left\{\left(1+\frac{1}{a_n}\right)^{a_n}\right\}$ 收敛于 $\mathrm{e}$ \\
(C) $\left\{a_n\ln a_n\right\}$ 收敛于 $0$ & (D) $\left\{\frac{\ln a_n}{a_n}\right\}$ 收敛于 $0$
\end{opts2}
% 自【解】起为解答，按规则截断不录。


% ==================== 第3~6讲（例3.1~例6.22，共48题） ====================
% ============ 第3讲（书页24~33, 36） ============

% ---- p01 ---- 书页24
\examp{例3.1}{设 $f(x)=2-x^{x}-x+o(x-1)$，$x\to 1$，且 $f'(1)=a$，则 $a=\underline{\hspace{3em}}$.}

% ---- p02~p03 ---- 书页25~26
\examp{例3.2}{设函数 $f(x)$ 在区间 $(-1,1)$ 内有定义，且在点 $x=0$ 处连续，则下列命题中：\\
(1) 当 $\lim\limits_{x\to 0}\dfrac{f(x)}{\sqrt[3]{x}}=0$ 时，$f(x)$ 在点 $x=0$ 处可导；\\
(2) 当 $\lim\limits_{x\to 0}\dfrac{f(x)}{x^{2}}=0$ 时，$f(x)$ 在点 $x=0$ 处可导；\\
(3) 当 $f(x)$ 在点 $x=0$ 处可导时，$\lim\limits_{x\to 0}\dfrac{f(x)}{\sqrt[3]{x}}=0$；\\
(4) 当 $f(x)$ 在点 $x=0$ 处可导时，$\lim\limits_{x\to 0}\dfrac{f(x)}{x^{2}}=0$.\\
真命题的个数为（\quad）.}
\begin{opts4}
(A) $1$ & (B) $2$ & (C) $3$ & (D) $4$
\end{opts4}

% ---- p04 ---- 书页28
\examp{例3.3}{设函数 $f(x)$ 连续，给出下列 $4$ 个条件：\\
(1) $\lim\limits_{x\to 0}\dfrac{|f(x)|-f(0)}{x}$ 存在；\\
(2) $\lim\limits_{x\to 0}\dfrac{f(x)-|f(0)|}{x}$ 存在；\\
(3) $\lim\limits_{x\to 0}\dfrac{|f(x)|}{x}$ 存在；\\
(4) $\lim\limits_{x\to 0}\dfrac{|f(x)|-|f(0)|}{x}$ 存在.\\
其中可得到“$f(x)$ 在 $x=0$ 处可导”的条件个数为（\quad）.}
\begin{opts4}
(A) $1$ & (B) $2$ & (C) $3$ & (D) $4$
\end{opts4}

% ---- p05~p06 ---- 书页29~30
\examp{例3.4}{设 $f(x)$ 在 $x=0$ 处连续，下列结论：\\
(1) 若 $f'_-(0)$ 存在，则 $f'_+(0)$ 存在；\\
(2) 若 $\lim\limits_{x\to 0^{-}}f'(x)$ 存在，则 $\lim\limits_{x\to 0^{+}}f'(x)$ 存在；\\
(3) 若 $f'_-(0)$ 与 $f'_+(0)$ 均存在，则 $f'(0)$ 存在；\\
(4) 若 $\lim\limits_{x\to 0}f'(x)$ 存在，则 $f'(0)$ 存在.\\
正确结论的个数为（\quad）.}
\begin{opts4}
(A) $1$ & (B) $2$ & (C) $3$ & (D) $4$
\end{opts4}

% ---- p07 ---- 书页30
\examp{例3.5}{设 $f(x)=x^{\frac{2}{3}}\sin x$，求 $f'(x)$.}

% ---- p08 ---- 书页31
\examp{例3.6}{设函数 $f(x)$ 处处可导，$f(0)=-1$，$f'(0)=1$，令 $g(x)=\bigl|f(x-1)\bigr|$，则 $g'(1)=\underline{\hspace{3em}}$.}

% ---- p09 ---- 书页32
\examp{例3.7}{下列函数中，在 $x=0$ 处不可导的是（\quad）.}
\begin{opts2}
(A) $f(x)=|x|\tan|x|$ & (B) $f(x)=|x|\tan\sqrt{|x|}$ \\
(C) $f(x)=\sqrt{\cos|x|}$ & (D) $f(x)=\cos\sqrt{|x|}$
\end{opts2}

% ---- p10 ---- 书页32续
\examp{例3.8}{已知函数 $g(x)$ 连续. 设 $f(x)=\int_{0}^{x^{2}}g(xt)\,\mathrm{d}t$，求 $f'(x)$ 的表达式，并判断 $f'(x)$ 在 $x=0$ 处的连续性.}

% ============ 第4讲（书页33~40） ============

% ---- p11 ---- 书页33
\examp{例4.1}{设函数 $f(x)$ 可导且满足 $x^{2}f'(x)=f^{2}(x)$，$f(1)=\dfrac{1}{3}$，则 $f^{(n)}(0)=$（\quad）.}
\begin{opts4}
(A) $(-1)^{n}n!$ & (B) $(-1)^{n-1}n!$ & (C) $(-2)^{n}n!$ & (D) $(-2)^{n-1}n!$
\end{opts4}

\examp{例4.2}{(1) 设 $y=\dfrac{1}{x(1-x)}$，求 $\dfrac{\mathrm{d}^{n}y}{\mathrm{d}x^{n}}$；(2) 设 $z=\dfrac{y^{2}}{x(1-x)}$，求 $\dfrac{\partial^{n}z}{\partial x^{n}}$.}

% ---- p12 ---- 书页36
\examp{例4.3}{设 $f(x)=(x^{3}-1)^{n}$，则 $f^{(n)}(1)=\underline{\hspace{3em}}$.}

\examp{例4.4}{设 $y=\dfrac{x}{1+x^{2}}$，求 $y^{(n)}$ 满足的递推关系式及 $y^{(2n+1)}(0)$.}

% ---- p13 ---- 书页37
\examp{例4.5}{设函数 $y=f(x)$ 由 $\begin{cases} x=2t+|t|,\\ y=|t|\tan t \end{cases}$ 所确定，则在 $\left(-\dfrac{\pi}{2},\dfrac{\pi}{2}\right)$ 内（\quad）.}
\begin{opts2}
(A) $f(x)$ 连续，$f'(0)$ 不存在 & (B) $f'(0)$ 存在，$f'(x)$ 在 $x=0$ 处不连续 \\
(C) $f'(x)$ 连续，$f''(0)$ 不存在 & (D) $f''(0)$ 存在，$f''(x)$ 在 $x=0$ 处不连续
\end{opts2}

% ---- p14 ---- 书页38
\examp{例4.6}{设 $f(x)=\dfrac{1}{2^{x}+1}$，$x\in\mathbf{R}$，则 $f^{(4)}(0)=\underline{\hspace{3em}}$.}

% ---- p15 ---- 书页39
\examp{例4.7}{已知可导函数 $y=y(x)$ 满足 $a\mathrm{e}^{x}+y^{2}+y-\ln(1+x)\cos y+b=0$，且 $y(0)=0$，$y'(0)=0$. (1) 求 $a,b$ 的值；(2) 判断 $x=0$ 是否为 $y(x)$ 的极值点.}

% ---- p16 ---- 书页39
\examp{例4.8}{若 $\begin{cases} x=\ln t,\\ y=\mathrm{e}^{t^{2}} \end{cases}$，则 $\left.\dfrac{\mathrm{d}^{2}y}{\mathrm{d}x^{2}}\right|_{t=1}=\underline{\hspace{3em}}$.}

% ---- p17 ---- 书页40
\examp{例4.9}{已知函数 $f(x)=\mathrm{e}^{x}+2x+1$，设 $g(y)$ 与 $f(x)$ 互为反函数，则 $g''(2)=$（\quad）.}
\begin{opts4}
(A) $\frac{1}{3}$ & (B) $-3$ & (C) $-\frac{1}{27}$ & (D) $-\frac{\mathrm{e}^{2}}{(\mathrm{e}^{2}+2)^{3}}$
\end{opts4}

% ---- p18 ---- 书页41
\examp{例5.1}{设 $f(x)$ 有连续的一阶导数，且 $f(0)=0$，$f'(0)=1$. 求极限 $\lim\limits_{x\to 0}\dfrac{xf(u)}{uf(x)}$，其中 $u$ 是曲线 $y=f(x)$ 在点 $(x,f(x))$ 处的切线在 $x$ 轴上的截距.}

% ============ 第5讲（书页41~47） ============

% ---- p19 ---- 书页43
\examp{例5.2}{设函数 $f(x)=(x^{2}+a)\mathrm{e}^{x}$，若 $f(x)$ 没有极值点，但曲线 $y=f(x)$ 有拐点，则 $a$ 的取值范围是（\quad）.}
\begin{opts4}
(A) $[0,1)$ & (B) $[1,+\infty)$ & (C) $[1,2)$ & (D) $[2,+\infty)$
\end{opts4}

\examp{例5.3}{（仅数学一、数学二）已知曲线 $y=f(x)$ 在点 $(0,1)$ 处的曲率圆方程为 $(x-1)^{2}+y^{2}=2$，且当 $x\to 0$ 时，二阶可导函数 $f(x)$ 与 $a+bx+cx^{2}$ 的差为 $o(x^{2})$，则（\quad）.}
\begin{opts2}
(A) $a=0$，$b=1$，$c=\frac{3}{2}$ & (B) $a=1$，$b=0$，$c=1$ \\
(C) $a=1$，$b=1$，$c=-1$ & (D) $a=1$，$b=0$，$c=-1$
\end{opts2}

% ---- p20~p21 ---- 书页44~45（例5.4 题干在44、选项在45，已合并）
\examp{例5.4}{设函数 $f(x)$ 在 $x=x_{0}$ 处有二阶导数，则（\quad）.}
\begin{opts1}
(A) 当 $f(x)$ 在 $x_{0}$ 的某邻域内单调增加时，$f'(x_{0})>0$ \\
(B) 当 $f'(x_{0})>0$ 时，$f(x)$ 在 $x_{0}$ 的某邻域内单调增加 \\
(C) 当曲线 $f(x)$ 在 $x_{0}$ 的某邻域内是凹的时，$f''(x_{0})>0$ \\
(D) 当 $f''(x_{0})>0$ 时，曲线 $f(x)$ 在 $x_{0}$ 的某邻域内是凹的
\end{opts1}

% ---- p22 ---- 书页45
\examp{例5.5}{曲线 $y=x\ln\left(\mathrm{e}+\dfrac{1}{x-1}\right)$ 的斜渐近线方程为（\quad）.}
\begin{opts4}
(A) $y=x+\mathrm{e}$ & (B) $y=x+\frac{1}{\mathrm{e}}$ & (C) $y=x$ & (D) $y=x-\frac{1}{\mathrm{e}}$
\end{opts4}

\examp{例5.6}{曲线 $y=\dfrac{(1+x)^{\frac{3}{2}}}{\sqrt{x}}$ 的斜渐近线方程为 $\underline{\hspace{3em}}$.}

% ---- p23 ---- 书页46
\examp{例5.7}{设 $g(x)$ 是函数 $f(x)=\dfrac{1}{2}\ln\dfrac{3+x}{3-x}$ 的反函数，则曲线 $y=g(x)$ 的渐近线方程为 $\underline{\hspace{3em}}$.}

\examp{例5.8}{$f(x)=\int_{0}^{x}\dfrac{t}{t^{2}+2t+2}\,\mathrm{d}t$ 在 $[0,1]$ 上的最大值为 $\underline{\hspace{3em}}$.}

% ---- p24 ---- 书页47
\examp{例5.9}{设 $f'(x)$ 在区间 $[0,4]$ 上连续，曲线 $y=f'(x)$ 与直线 $x=0$，$x=4$，$y=0$ 围成如图所示的三个区域，其面积分别为 $S_{1}=3$，$S_{2}=4$，$S_{3}=2$，且 $f(0)=1$，则 $f(x)$ 在 $[0,4]$ 上的最大值与最小值分别为（\quad）.}
\begin{opts4}
(A) $2,-3$ & (B) $4,-3$ & (C) $2,-2$ & (D) $4,-2$
\end{opts4}

% ============ 第6讲（书页49~60） ============

% ---- p25 ---- 书页49
\examp{例6.1}{设函数 $f(x)$ 在 $[0,1]$ 上连续，$(0,1)$ 内可导，$f(0)=f(1)$，$\int_0^1 f(x)\,\mathrm{d}x=0$，证明：存在 $\xi\in(0,1)$，使得 $f'(\xi)+f^2(\xi)=0$.}

% ---- p26 ---- 书页52
\examp{例6.2}{已知函数 $f(x)$ 和 $g(x)$ 在 $[a,b]$ 上连续，在 $(a,b)$ 内可导，且 $g'(x)\neq 0$，证明：存在 $\xi\in(a,b)$，使得 $\dfrac{f(\xi)-f(a)}{g(b)-g(\xi)}=\dfrac{f'(\xi)}{g'(\xi)}$.}

\examp{例6.3}{设 $f(x)$ 在 $[0,1]$ 上连续，在 $(0,1)$ 内可导，且满足 $f(1)=k\int_0^{\frac{1}{k}}x\mathrm{e}^{1-x}f(x)\,\mathrm{d}x\ (k>1)$，证明至少存在一点 $\xi\in(0,1)$，使得 $f'(\xi)=\left(1-\xi^{-1}\right)f(\xi)$.}

% ---- p27 ---- 书页52~53
\examp{例6.4}{设正值函数 $f(x)$，$g(x)$ 在 $[a,b]$ 上连续，证明存在 $\xi\in(a,b)$，使得 $\dfrac{f(\xi)}{g(\xi)}=\dfrac{\displaystyle\int_a^\xi f(x)\,\mathrm{d}x}{\displaystyle\int_\xi^b g(x)\,\mathrm{d}x}$.}

% ---- p28 ---- 书页53
\examp{例6.5}{证明：若 $f(x)$ 连续且满足 $\int_0^{\frac{\pi}{2}}f(x)\cos x\,\mathrm{d}x=0$，则存在 $\xi\in\left(0,\frac{\pi}{2}\right)$，使得 $f(\xi)=0$.}

% ---- p29 ---- 书页54
\examp{例6.6}{设函数 $f(x)$ 在 $[0,\pi]$ 上连续，且 $\int_0^\pi f(x)\,\mathrm{d}x=0$，$\int_0^\pi f(x)\cos x\,\mathrm{d}x=0$. 证明：在 $(0,\pi)$ 内至少存在两个不同的点 $\xi_1$，$\xi_2$，使得 $f(\xi_1)=f(\xi_2)=0$.}

% ---- p30 ---- 书页55
\examp{例6.7}{设函数 $f(x)$ 在 $[0,1]$ 上二阶可导，$f(0)=0$，且 $f(x)$ 在 $(0,1)$ 内取得最大值 $2$，在 $(0,1)$ 内取得最小值，证明：(1) 存在 $\xi\in(0,1)$，使得 $f'(\xi)>2$；(2) 存在 $\eta\in(0,1)$，使得 $f''(\eta)<-4$.}

% ---- p31 ---- 书页55
\examp{例6.8}{(1) 设 $f(x)$ 在 $[a,b]$ 上可导，若 $f'_{+}(a)\neq f'_{-}(b)$，证明：对于任意的介于 $f'_{+}(a)$ 与 $f'_{-}(b)$ 之间的 $\mu$，存在 $\xi\in(a,b)$，使得 $f'(\xi)=\mu$；\\
(2) 若 $f(x)$ 在 $[0,2]$ 上具有二阶导数，证明：存在 $\xi\in(0,2)$，使得 $f(0)+f(2)-2f(1)=f''(\xi)$.}

% ---- p32 ---- 书页56
\examp{例6.9}{已知函数 $f(x)$ 在 $[x_{0},\,x_{0}+\delta)$ 上连续，在 $(x_{0},\,x_{0}+\delta)$ 内可导，$\delta>0$，证明：若 $\lim\limits_{x\to x_{0}^{+}}f'(x)=A$，则 $f'_{+}(x_{0})=A$.}

% ---- p33 ---- 书页57
\examp{例6.10}{设函数 $f(x)$ 在区间 $[a,b]$ 上满足：对任意 $x$, $y\in[a,b]$，有 $|f(x)-f(y)|\leqslant M|x-y|^{\alpha}$，其中 $M>0$，$\alpha>1$ 是常数. 证明：$f(x)$ 在 $[a,b]$ 上恒为常数.}

% ---- p34 ---- 书页58
\examp{例6.11}{已知函数 $f(x)$ 在 $(-\infty,0)$ 上可导，且 $\lim\limits_{x\to-\infty}f'(x)=A>0$，证明 $\lim\limits_{x\to-\infty}f(x)=-\infty$.}

\examp{例6.12}{设函数 $f(x)$ 在区间 $[0,2]$ 上具有连续导数，$f(0)=f(2)=0$，$M=\max\limits_{x\in[0,2]}\{|f(x)|\}$. 证明：(1) 存在 $\xi\in(0,2)$，使得 $|f'(\xi)|\geqslant M$；(2) 若对任意的 $x\in(0,2)$，$|f'(x)|\leqslant M$，则 $M=0$.}

% ---- p35 ---- 书页59
\examp{例6.13}{已知 $\int_{0}^{x}f(t)\,\mathrm{d}t=xf(\theta x)$，$x>0$，$f(x)=\mathrm{e}^{x}$，则 $\lim\limits_{x\to 0^{+}}\theta=\underline{\hspace{3em}}$.}

% ---- p36 ---- 书页60
\examp{例6.14}{设 $f(x)$ 在 $[a,b]$ 上连续，且 $f(x)>0$，证明存在 $\xi\in(a,b)$，使得 $\dfrac{b^{2}-a^{2}}{\displaystyle\int_{a}^{b}f(x)\,\mathrm{d}x}=\dfrac{2\xi}{f(\xi)}$.}

% ---- p37 ---- 第6讲综合题举例
\examp{例6.15}{设 $f(x)=\int_1^{x}\sin t^2\,\mathrm{d}t$，$g(x)=\int_0^{x}f(t)\,\mathrm{d}t$. (1) 计算 $g(1)$；(2) 证明：存在 $\xi\in(0,1)$，使得 $\int_1^{\xi}\sin t^2\,\mathrm{d}t=\xi(\cos 1-1)$.}

% ---- p38 ----
\examp{例6.16}{设函数 $f(x)$ 在 $[0,1]$ 上二阶可导，$f(0)=f(1)$，且 $|f''(x)|\leqslant 2$，证明：$|f'(x)|\leqslant 1$，$x\in[0,1]$.}

% ---- p39 ----
\examp{例6.17}{设函数 $f(x)=ax-b\ln x$（$a>0$）有 $2$ 个零点，则 $\dfrac{b}{a}$ 的取值范围是（\quad）.}
\begin{opts4}
(A) $(0,\mathrm{e})$ & (B) $(\mathrm{e},+\infty)$ & (C) $\left(0,\frac{1}{\mathrm{e}}\right)$ & (D) $\left(\frac{1}{\mathrm{e}},+\infty\right)$
\end{opts4}

\examp{例6.18}{设 $f(x)$ 在 $\left[0,\dfrac{3}{2}\pi\right]$ 上连续，在 $\left(0,\dfrac{3}{2}\pi\right)$ 内是函数 $\dfrac{\sin x}{x}$ 的一个原函数，$f(0)=0$. (1) 证明 $f\left(\dfrac{3}{2}\pi\right)>0$；(2) 求方程 $\int_0^{x}\dfrac{\sin t}{t}\,\mathrm{d}t=\ln x^2$ 的实根个数.}

% ---- p40 ----
\examp{例6.19}{设函数 $f(x)$ 在区间 $(a,b)$ 内可导. 证明：导函数 $f'(x)$ 在 $(a,b)$ 内严格单调增加的充分必要条件是对 $(a,b)$ 内任意的 $x_1$，$x_2$，$x_3$，当 $x_1<x_2<x_3$ 时，$\dfrac{f(x_2)-f(x_1)}{x_2-x_1}<\dfrac{f(x_3)-f(x_2)}{x_3-x_2}$.}

% ---- p41 ----
\examp{例6.20}{已知常数 $k\geqslant\ln 2-1$. 证明：$(x-1)\left(x-\ln^2 x+2k\ln x-1\right)\geqslant 0$.}

\examp{例6.21}{设函数 $y=y(x)$ 由参数方程 $\begin{cases} x=\dfrac{1}{3}t^3+t+\dfrac{1}{3},\\ y=\dfrac{1}{3}t^3-t+\dfrac{1}{3} \end{cases}$ 确定. \\
(1) 求 $y(x)$ 的极值；\\
(2) 若 $\begin{cases} u=x,\\ v=y+k, \end{cases}$ 且 $v=v(u)$ 恰有一个零点，求常数 $k$ 的取值范围.}

% ---- p42 ----
\examp{例6.22}{设可导函数 $f(x)$ 严格单调递增且满足 $\int_{-1}^{1}f(x)\,\mathrm{d}x=0$，记 $a=\int_0^{1}f(x)\,\mathrm{d}x$. \\
(1) 证明 $a>0$；\\
(2) 令 $F(x)=a(1-x^2)+\int_1^{x}f(t)\,\mathrm{d}t$，证明：存在 $\xi\in(-1,1)$，使得 $F''(\xi)=0$.}

\end{document}
